\section*{Bab \ref{ch:g11:infrared} --- Spektroskopi Inframerah}

\begin{solution}{exo:g11:infrared:1}
$\lambda = 1/3300 = \qty{3.03e-4}{cm} = \qty{3.03}{\micro m}$;
$\lambda = 1/1050 = \qty{9.52e-4}{cm} = \qty{9.52}{\micro m}$. Pita pada
\qty{3300}{cm^{-1}}, yang \omterm{def:g11:infrared:wavenumber}{bilangan gelombangnya} lebih besar, milik
getaran yang lebih cepat.
\end{solution}

\begin{solution}{exo:g11:infrared:2}
$\qty{10}{\micro m} = \qty{1.0e-3}{cm}$, jadi $\sigma = 1/\num{1.0e-3} =
\qty{1000}{cm^{-1}}$.
\end{solution}

\begin{solution}{exo:g11:infrared:3}
\ce{C=O} \omterm{def:g11:functional-groups:carbonyl}{keton}; \ce{O-H} \omterm{def:g11:functional-groups:alcohol}{alkohol} yang berikatan hidrogen; ikatan
\ce{C-H} rantai karbon.
\end{solution}

\begin{solution}{exo:g11:infrared:4}
Pita pada \qty{1715}{cm^{-1}}: ikatan \ce{C=O}.
\end{solution}

\begin{solution}{exo:g11:infrared:5}
Alat itu mengukur cahaya yang menembus sampel; di tempat sebuah ikatan
menyerap, cahaya yang lolos lebih sedikit dan kurvanya melekuk ke bawah.
$T =
\qty{100}{\percent}$ berarti tidak ada cahaya dengan \omterm{def:g11:infrared:wavenumber}{bilangan
gelombang} itu yang diserap.
\end{solution}

\begin{solution}{exo:g11:infrared:6}
Yang pertama memiliki pita \ce{C=O} pada \qty{1715}{cm^{-1}} dan tidak
memiliki pita \ce{O-H}. Senyawa itu \omterm{def:g11:functional-groups:carbonyl}{keton}: \ce{CH3-CO-CH2-CH3}, butanon
(juga disebut butan-2-on). Yang kedua memiliki pita \ce{O-H} dan tidak
memiliki \ce{C=O}: sebuah \omterm{def:g11:functional-groups:alcohol}{alkohol}, misalnya butan-1-ol.
\end{solution}

\begin{solution}{exo:g11:infrared:7}
Gugus-gugus \ce{O-H} pada asam itu berikatan hidrogen dengan sangat kuat
(\omterm{def:g7:atoms-and-molecules:molecule}{molekul-molekulnya} berpasangan, setiap \ce{O-H} berikatan dengan
\ce{C=O} \omterm{def:g7:atoms-and-molecules:molecule}{molekul} pasangannya), dan setiap \omterm{def:g7:atoms-and-molecules:molecule}{molekul} terikat sedikit
berbeda: pitanya makin dilonggarkan dan tersebar pada rentang yang
sangat lebar.
\end{solution}

\begin{solution}{exo:g11:infrared:8}
Dalam keadaan cair (A), pita \ce{O-H}-nya lebar, di sekitar
\qty{3350}{cm^{-1}}; dalam keadaan gas (B), tempat \omterm{def:g7:atoms-and-molecules:molecule}{molekul-molekulnya}
berjauhan dan tidak membentuk \omterm{def:g11:polarity-and-cohesion:hydrogen-bond}{ikatan hidrogen}, pita itu tajam di sekitar
\qty{3666}{cm^{-1}}.
\end{solution}

\begin{solution}{exo:g11:infrared:9}
Hampir tidak: sekitar \qty{1730}{cm^{-1}} untuk \omterm{def:g11:functional-groups:carbonyl}{aldehida} dan
\qty{1715}{cm^{-1}} untuk \omterm{def:g11:functional-groups:carbonyl}{keton}, nilai yang berdekatan. Membandingkan
seluruh spektrum, termasuk \omterm{def:g11:infrared:band}{daerah sidik jarinya}, dengan spektrum acuan
akan memutuskannya, atau dengan teknik lain.
\end{solution}

\begin{solution}{exo:g11:infrared:10}
Sebuah \omterm{def:g11:functional-groups:carboxylic-acid}{ester} (\ce{C=O} di sekitar \qty{1735}{cm^{-1}}, \ce{C-O} yang
kuat, tanpa \ce{O-H}), misalnya etil etanoat. Asam etanoat adalah
\ce{C2H4O2}, bukan \ce{C4H8O2}. Asam butanoat memiliki rumus yang tepat
tetapi akan memperlihatkan pita \ce{O-H} asam yang sangat lebar:
dikesampingkan.
\end{solution}

\begin{solution}{exo:g11:infrared:11}
Propan-1-ol adalah \omterm{def:g11:functional-groups:alcohol}{alkohol}: seperti A. Asam propanoat adalah \omterm{def:g11:functional-groups:carboxylic-acid}{asam
karboksilat}: seperti D.
\end{solution}

\begin{solution}{exo:g11:infrared:12}
Pita \ce{O-H} yang lebar di sekitar \qty{3350}{cm^{-1}} mengecil
sementara pita \ce{C=O} yang kuat di sekitar \qty{1715}{cm^{-1}}
membesar. Reaksinya selesai ketika pita \ce{O-H} telah hilang dan pita
\ce{C=O} tidak membesar lagi.
\end{solution}

\begin{solution}{exo:g11:infrared:13}
Asam butanoat memperlihatkan pita \ce{O-H} yang sangat lebar dari 2500
sampai \qty{3300}{cm^{-1}}, \omterm{def:g11:functional-groups:carboxylic-acid}{esternya} tidak; \ce{C=O} asam terletak di
sekitar \qty{1710}{cm^{-1}}, lebih lebar, sedangkan milik \omterm{def:g11:functional-groups:carboxylic-acid}{ester} di sekitar
\qty{1735}{cm^{-1}}, tajam, dengan pita \ce{C-O} yang kuat di sekitar
\qty{1240}{cm^{-1}}. Asam memiliki gugus \ce{O-H} yang berikatan
hidrogen, yang juga sedikit melonggarkan \ce{C=O}-nya.
\end{solution}

\begin{solution}{exo:g11:infrared:14}
Sisa etanol (\omterm{def:g11:functional-groups:alcohol}{alkohol} bahan pembuat \omterm{def:g11:functional-groups:carboxylic-acid}{ester} itu) atau air (dari \omterm{def:g4:air-a-mixture-of-gases:air}{udara} atau
dari pencucian): keduanya memiliki gugus \ce{O-H} yang berikatan
hidrogen. Sampel itu dapat dibandingkan dengan spektrum acuan,
dikeringkan, didistilasi, atau diukur suhu didihnya.
\end{solution}

\begin{solution}{exo:g11:infrared:15}
\omterm{def:g10:lewis-and-shape:multiple-bond}{Ikatan rangkap dua} adalah pegas yang lebih kaku daripada ikatan tunggal
antara \omterm{def:g7:atoms-and-molecules:atom}{atom-atom} yang sama: ikatan itu bergetar lebih cepat, pada
\omterm{def:g11:infrared:wavenumber}{bilangan gelombang} yang lebih tinggi (1715 dibandingkan
\qty{1050}{cm^{-1}}). Ikatan \ce{O-H}, \ce{N-H}, dan \ce{C-H} semuanya
melibatkan \omterm{def:g7:atoms-and-molecules:atom}{atom} hidrogen, yang paling ringan: bola ringan pada pegas
bergetar cepat, sehingga \omterm{def:g11:infrared:wavenumber}{bilangan gelombangnya} sekitar
2900--\qty{3600}{cm^{-1}}.
\end{solution}

\begin{solution}{pb:g11:infrared:1}
\textbf{1.} \ce{CH3-CH2-OH}; \ce{CH3-CO-CH3}; \ce{CH3-COOH}.

\textbf{2.} \ce{C2H6O}; \ce{C3H6O}; \ce{C2H4O2}.

\textbf{3.} Hidroksil; karbonil (\omterm{def:g11:functional-groups:carbonyl}{keton}); karboksil.

\textbf{4.} Etanol: \ce{O-H} dan \ce{C-H}. Propanon: \ce{C-H} dan
\ce{C=O}. Asam etanoat: \ce{O-H}, \ce{C-H}, dan \ce{C=O}.

\textbf{5.} Tidak ada pita \ce{C=O}, ada pita \ce{O-H} yang lebar:
etanol.

\textbf{6.} Pita \ce{O-H} yang sangat lebar dari 2500 sampai
\qty{3300}{cm^{-1}} dan \ce{C=O} yang kuat di sekitar
\qty{1710}{cm^{-1}}: asam etanoat.

\textbf{7.} Propanon: pita \ce{C=O} yang sangat kuat pada
\qty{1715}{cm^{-1}}, tanpa \ce{O-H}.

\textbf{8.} Mungkin (cuka, \omterm{def:g6:solutions-and-solubility:solute}{pelarut}, \omterm{def:g11:functional-groups:alcohol}{alkohol}), tetapi mengendus cairan
yang tidak diketahui itu berbahaya; spektrum aman dan pasti.

\textbf{9.} Gugus-gugus \ce{O-H} asam berikatan hidrogen lebih kuat
(\omterm{def:g7:atoms-and-molecules:molecule}{molekul-molekulnya} berpasangan), sehingga pitanya lebih rendah dan
jauh lebih lebar.

\textbf{10.} \ce{CH3-O-CH3}, \ce{C2H6O}; senyawa itu tidak memiliki
\ce{O-H}, jadi tidak ada pita \ce{O-H}.

\textbf{11.} Pita \ce{C=O}: sekitar \qty{1730}{cm^{-1}} untuk propanal,
sekitar \qty{1715}{cm^{-1}} untuk propanon.

\textbf{12.} Pita \ce{O-H}: \omterm{def:g11:functional-groups:carboxylic-acid}{ester} tidak memiliki \ce{O-H}. Pita
\ce{C=O}-nya akan terletak di sekitar \qty{1735}{cm^{-1}}.

\textbf{13.} \omterm{def:g11:organic-skeletons:isomer}{Isomer-isomer} memiliki rumus yang sama; \omterm{def:g11:functional-groups:functional-group}{gugus fungsinya}
berbeda, dan spektrum memperlihatkan gugus-gugus itu.

\textbf{14.} \qty{1715}{cm^{-1}}.

\textbf{15.} $1715 \times 100 = \qty{1.715e5}{m^{-1}}$.

\textbf{16.} $\lambda = 1/\num{1.715e5} = \qty{5.83e-6}{m} =
\qty{5.83}{\micro m}$.

\textbf{17.} Tidak: cahaya tampak terbentang dari sekitar 0.4 sampai
\qty{0.75}{\micro m}. Radiasi itu adalah radiasi inframerah.

\textbf{18.} $\num{1.05e5}\ \si{m^{-1}}$, jadi $\lambda = \qty{9.52e-6}{m}
= \qty{9.52}{\micro m}$.

\textbf{19.} Ikatan \ce{C=O} (\omterm{def:g11:infrared:wavenumber}{bilangan gelombang} lebih tinggi, panjang
gelombang lebih pendek).

\textbf{20.} Sekitar \qty{5.8}{\micro m}.
\end{solution}
